% =====================================================================
%  DRAFT -- Re-embed, Map, or Window?
%  Characterizing Reuse Strategies for Embedding Model Upgrades
%
%  Status: skeleton + method + experimental protocol.
%  NO RESULTS ARE REAL. Every number to be filled is marked \todo{}.
%  Target: Experiments, Analysis & Benchmarks track (ICDE / VLDB / SIGMOD).
%  Anonymity: Thaw is cited as anonymous concurrent work; adjust to venue rules.
% =====================================================================
\documentclass[conference]{IEEEtran}
\IEEEoverridecommandlockouts
\usepackage{amsmath,amssymb,amsfonts}
\usepackage{algorithmic}
\usepackage{graphicx}
\usepackage{booktabs}
\usepackage{multirow}
\usepackage{xcolor}
\usepackage{url}
\usepackage[hidelinks]{hyperref}

\newcommand{\todo}[1]{\textcolor{red}{[TODO: #1]}}
\newcommand{\note}[1]{\textcolor{blue}{[NOTE: #1]}}
\newcommand{\topk}{\mathrm{Top}}
\newcommand{\win}{W}
\newcommand{\Gk}{G(K')}

\begin{document}

\title{Re-embed, Map, or Window? Characterizing Reuse Strategies\\for Embedding Model Upgrades}

\author{\IEEEauthorblockN{Anonymous Authors}
\IEEEauthorblockA{\textit{Under double-blind review}}}

\maketitle

% ---------------------------------------------------------------------
\begin{abstract}
Upgrading the embedding model behind a vector index normally means re-encoding the
whole corpus before the new model can serve. Several reuse strategies avoid or defer
that cost: mapping new-model queries into the old space with a learned adapter,
re-ranking a window of candidates nominated by the old index with the new model, and
combinations of these with model-independent lexical nomination and background
migration. Each has been evaluated in isolation, on its own metric, and mostly on
text. We do not propose another strategy. Instead, we evaluate all of them under a
single cost model --- recall along the path to the \emph{target} model as a function of
encoding work --- across \todo{N} model pairs spanning same-family, cross-family,
quantization, and code-specialized upgrades, on text and code corpora.
We ask three questions: which strategy wins where; whether two cheap pair-level
measurements (space compatibility and window residual miss) predict the winner on a
held-out model pair; and whether that prediction can be made \emph{before} migration
at bounded encoding cost. We additionally test, under a pre-registered gate, whether a
document-level signal can order background migration better than random.
\todo{One-sentence headline findings once the regime map is measured.}
All code, encodings, and per-pair measurements are released.
\end{abstract}

\begin{IEEEkeywords}
vector search, embedding models, index maintenance, model migration, benchmark
\end{IEEEkeywords}

% =====================================================================
\section{Introduction}
% =====================================================================

Dense retrieval systems tie their index to a specific embedding model. When a better
model appears, the vectors in the index become incompatible with the new query
encoder, and the standard practice is a blue--green rebuild: re-encode the corpus,
build a second index, validate it, and switch traffic~\cite{zilliz-migration}. At the
scale of hundreds of millions of documents this costs days to weeks of accelerator
time before the new model serves a single query.

Three families of reuse strategies have emerged to avoid that up-front cost.
\emph{Space mapping} learns a light transformation that maps new-model queries into the
old space so the existing index keeps serving~\cite{drift-adapter}.
\emph{Window re-ranking} retrieves a bounded candidate window from the old index,
scores it with the new model, and materializes new vectors progressively in the
background~\cite{embedflow,thaw-anon}. \emph{Structure reuse} exploits the observation
that old-model neighborhoods remain partially informative under the new model, for
instance to accelerate building the new graph index~\cite{dgm}.

These strategies were evaluated on different metrics, different model pairs, and
almost exclusively on text. A practitioner facing a concrete upgrade --- say a
general-purpose encoder to a code-specialized one --- has no basis for choosing among
them, and the published numbers cannot be compared. Moreover, space mapping does not
upgrade stored vectors: it bounds retrieval by what the old model encoded, so it is a
bridge, not a destination. Any fair comparison must therefore measure the cost of
\emph{reaching the new model} while keeping recall acceptable along the way, not only
the quality of a frozen intermediate state.

\textbf{Research questions.}
\begin{itemize}
\item \textbf{RQ1 (Regime map).} For each model pair, which reuse strategy achieves the
best recall-versus-encoding-cost trajectory toward the target model, and how does the
winner vary across pair types and corpora?
\item \textbf{RQ2 (Predictability).} Do two pair-level measurements --- space
compatibility and window residual miss --- predict the winning strategy on a model
pair that was not used to build the map?
\item \textbf{RQ3 (Bounded-cost probe).} Can those measurements be estimated before
migration while encoding only a small, bounded fraction of the corpus, and how large
is the estimation error relative to the full offline measurement?
\item \textbf{RQ4 (Document-level ordering, gated).} Within a chosen strategy, does a
pre-embedding document signal order background migration better than random, size,
or demand-only order at equal encoding budget?
\end{itemize}

\textbf{Contributions.}
\begin{enumerate}
\item \textbf{A unified evaluation frame} that scores every reuse strategy by its
recall trajectory toward the target model per unit of encoding work, compatible with
the window gap $G(K')$ of~\cite{embedflow} (Section~\ref{sec:frame}).
\item \textbf{A measured regime map} over \todo{N} model pairs on text and code, with
a falsifiable claim tested by leave-one-pair-out winner prediction
(Sections~\ref{sec:rq1}--\ref{sec:rq2}).
\item \textbf{A bounded-cost probe} that estimates window residual miss before
migration via an intruder-count estimator, validated against full offline
measurements (Section~\ref{sec:rq3}).
\item \textbf{A pre-registered test of document-level migration ordering}, reported
whether positive or negative (Section~\ref{sec:rq4}).
\end{enumerate}

\textbf{Scope.} We study which existing strategy wins where and how to tell cheaply in
advance. We make no exactness guarantees for mixed-state retrieval and propose no new
serving mechanism; those questions are treated in concurrent
work~\cite{thaw-anon}.

\todo{Summary of findings paragraph (3--4 bullets) after RQ1--RQ4 are run.}

% =====================================================================
\section{Problem Setting and Evaluation Frame}\label{sec:frame}
% =====================================================================

\subsection{Setting}
Let $D$ be a corpus of $N$ documents, $f_s$ the source (old) encoder, and $f_t$ the
target (new) encoder. An index $I_s$ over $\{f_s(d)\}_{d\in D}$ is serving. A query
workload $Q$ arrives over time, possibly with repeated queries. Let
$\topk_K^{t}(q)$ denote the native top-$K$ of $q$ under $f_t$ over all of $D$, and
$\topk_{K'}^{s}(q)$ the top-$K'$ window returned by $I_s$.

\subsection{Cost model}
The dominant cost of migration is encoding documents with $f_t$. We count
\emph{document encodings} $C$ as the primary cost unit and report wall-clock and
accelerator time as secondary units. Query encoding, ANN search, and adapter
application are measured but are orders of magnitude cheaper per call;
\todo{confirm with measurements and report the ratio}.

\subsection{Metrics}

\textbf{Reference qualities.} Let $M_T$ be retrieval quality (Recall@10 or nDCG@10
against qrels) of exact native $f_t$ retrieval over all of $D$, and $M_S$ the same for
the source model $f_s$ --- the quality the system serves before migration starts. For a
strategy $\sigma$ after $C$ encodings, $M_\sigma(C)$ is the quality it serves at that point.

\textbf{Window gap.} Following~\cite{embedflow}, for window depth $K'$,
\[
G(K') = M_T - M_{T\mid S_{K'}},
\]
where $M_{T\mid S_{K'}}$ is the quality of $f_t$ re-ranking the source window only.

\textbf{Residual miss.} The qrels-free counterpart, measured against the native target
ranking:
\[
\rho(K') = 1 - \mathbb{E}_{q\in Q}\,
\frac{\lvert \topk_K^{t}(q) \cap \topk_{K'}^{s}(q)\rvert}{K}.
\]
Because $\rho$ depends on $K'/N$, we always report $K'$ as a fraction of $N$.
All windows of different depths are prefixes of a single retrieval at the largest depth,
so $\rho(K')$ is non-increasing in $K'$ by construction; measurements use exact search,
and the approximate-search noise floor is reported separately (pair P0).

\textbf{Trajectory metrics.} For each strategy we report (i) the curve
$C \mapsto M_\sigma(C)$ from $C{=}0$ to full migration; (ii) the \emph{cost-to-$\varepsilon$},
$C_\varepsilon(\sigma) = \min\{C : M_\sigma(C) \ge M_T - \varepsilon\}$; and (iii) the
\emph{floor} $\min_C M_\sigma(C)$, the worst quality served during migration.
Space mapping alone never reaches $M_T$; it is evaluated as a bridge combined with a
background rebuild.

\textbf{Pre-registered winner rule.} A strategy--order combination is \emph{admissible}
on a pair if it never serves worse than the system it replaces,
$\min_C M_\sigma(C) \ge M_S - \delta$. The \emph{winner} is the admissible combination
with the smallest $C_{\varepsilon}$, ties broken by the area under the trajectory. We fix
$\varepsilon = 0.01$ and $\delta = \todo{0.01}$ on MS MARCO dev Recall@10 before any
pair is run; nDCG@10 on TREC-DL and the other trajectory metrics are secondary. The
floor is anchored to $M_S$, not $M_T$: at $C{=}0$ every strategy serves roughly $M_S$,
so a target-anchored floor would exclude all strategies whenever the upgrade is worth doing.

\textbf{Pair-level measurements.} \emph{Compatibility} $\kappa = 1 - r$, where $r$ is the
relative residual $\lVert A(X_t) - X_s\rVert_F / \lVert X_s\rVert_F$ of a map $A$ from $f_t$ to
$f_s$ embeddings (semi-orthogonal Procrustes, or ridge-regularized affine), fit on one
document sample and evaluated on a disjoint one; larger $\kappa$ means more compatible.
We also report linear CKA~\cite{cka} on the same sample.
\emph{Window reliability}: $\rho(K')$ and $G(K')$ at fixed $K'/N$.

% =====================================================================
\section{Strategies Under Study}\label{sec:strategies}
% =====================================================================

Table~\ref{tab:strategies} lists the strategies. Every strategy that does not reach
the target model on its own is paired with a background migration that re-encodes the
corpus at a fixed rate, so all trajectories end at $M_T$.

\begin{table*}[t]
\centering
\caption{Strategies evaluated. ``Nominator'' is what proposes candidates; ``Ranker''
is what orders them.}
\label{tab:strategies}
\small
\begin{tabular}{@{}llll@{}}
\toprule
ID & Strategy & Nominator & Ranker \\
\midrule
S0 & Blue--green rebuild & --- (old index until switch) & $f_s$, then $f_t$ \\
S1 & Space mapping~\cite{drift-adapter} & $I_s$ via $A(f_t(q))$ & $f_s$ \\
S2 & Window re-rank~\cite{embedflow,thaw-anon} & $I_s$ via $f_s(q)$ & $f_t$ \\
S3 & Mapped window & $I_s$ via $A(f_t(q))$ & $f_t$ \\
S4 & Window $\cup$ lexical & $I_s \cup$ BM25~\cite{bm25} & $f_t$ \\
S5 & Window $\cup$ mapped $\cup$ lexical & all three & $f_t$ \\
\bottomrule
\end{tabular}
\end{table*}

\textbf{New writes.} Documents inserted during migration have no source vector and can
never be nominated by $I_s$. All window strategies therefore also search a small
target-space index over migrated and newly written documents and merge results.

\textbf{Background migration orders.} For S2--S5 the background re-encoder follows one
of: \emph{O-rand} (uniform random), \emph{O-size} (shortest documents first, maximizing
documents per GPU-second), \emph{O-demand} (no background work: only documents surfaced
by queries are encoded, i.e.\ pure lazy materialization, which need not ever complete),
and \emph{O-signal} (Section~\ref{sec:rq4}). Window documents missing a target vector are
either encoded synchronously when a query is served (\emph{sync}) or queued ahead of
background work (\emph{async}, as in~\cite{embedflow}).

\textbf{What a checkpoint measures.} Quality along the trajectory is measured from the
migration state, not from hypothetical encodings. A query whose window is fully encoded
(\emph{warm}) is ranked by target scores over its window and the target index; a query
with any unencoded window document (\emph{cold}) is served its nomination-order answer.
Scoring unencoded window documents at a checkpoint would let window strategies reach
target quality at $C{=}0$ without paying for those encodings, and mixing target-scored
with unscored documents in one ranking compares scores across spaces. For sync strategies
we separately report the number of encodings performed per served query, a latency proxy.

% =====================================================================
\section{Experimental Setup}\label{sec:setup}
% =====================================================================

\subsection{Corpora and queries}
\begin{itemize}
\item \textbf{Text:} MS MARCO passages~\cite{msmarco}. The full-encoding experiments use
a \todo{1M}-passage subsample built as \emph{every judged passage} (MS MARCO dev and
TREC-DL 2019/2020 qrels) plus a uniform random fill; a random subsample would drop most
judged passages and make $M_T$ and $G(K')$ meaningless. Recall@10 is computed on MS MARCO
dev (\todo{$\approx$7k} queries, sparse labels; primary) and nDCG@10 on TREC-DL 2019+2020
(\todo{$\approx$100} queries, dense labels; secondary). Full-collection runs provide scale points.
\item \textbf{Text with real heat:} the MS MARCO document collection with ORCAS click
logs~\cite{orcas}, used to construct query workloads with empirical repetition.
\note{ORCAS clicks are document-level; map to passages if passages are used.}
\item \textbf{Code:} CodeSearchNet~\cite{codesearchnet} (Python, Go, JavaScript) and the
repositories of CIPBench~\cite{cipbench} as small real-evolution sanity checks.
\end{itemize}

\textbf{Workloads.} Uniform over the query set; ORCAS empirical frequency (text); and a
Zipf$(\alpha)$ resampling with $\alpha\in\{0, 0.6, 0.9, 1.2\}$ (code), so that every
skew-dependent claim is reported as a function of $\alpha$ rather than assumed.

\textbf{Query roles.} A seeded fraction \todo{0.5} of the judged queries is held out for
evaluation and TREC-DL is evaluation-only; all other queries form the traffic that is
served during migration. Evaluated queries enter the traffic only as explicit repeats, at
rate $r\in$\todo{$\{0, 0.3\}$}. Without this split, on-demand repair warms exactly the
windows being evaluated and wins by construction; an earlier version of our replay had
this flaw. \textbf{Traffic intensity.} With background budget $b = N/R$ encodings per round,
each round serves $\lambda$ queries, set by the ratio
$\tau = \lambda K'/b \in$\todo{$\{0.25, 1, 4\}$}: the on-demand encodings a fully cold
round could trigger, relative to the background budget. We report the winner at every
$(\tau, r)$ and pre-register one setting, \todo{$(\tau{=}1, r{=}0)$}, as the label RQ2 predicts.

\subsection{Model pairs}
Table~\ref{tab:pairs} lists the planned pairs. P0 is the control: the same model
with its approximate index rebuilt under a different seed, so its $\rho(K')$ measures
pure ANN noise and serves as the noise floor for every other pair. P6 (fp32$\to$int8)
is only a sanity check: it transforms stored vectors and needs no re-encoding, so every
strategy wins trivially and it is excluded from winner prediction.

\begin{table*}[t]
\centering
\caption{Planned model pairs (source $\rightarrow$ target). \todo{verify dims and
availability; add pairs to reach 8--10.}}
\label{tab:pairs}
\small
\begin{tabular}{@{}lllc@{}}
\toprule
ID & Source & Target & Type \\
\midrule
P0 & bge-base-en-v1.5 & same, ANN rebuilt & control \\
P1 & all-MiniLM-L6-v2 & all-mpnet-base-v2 & cross-arch \\
P2 & all-MiniLM-L6-v2 & bge-base-en-v1.5 & cross-family \\
P3 & bge-small-en-v1.5 & bge-base-en-v1.5 & same family \\
P4 & e5-base-v2 & bge-base-en-v1.5 & cross-family \\
P5 & Qwen3-Emb-0.6B & Qwen3-Emb-4B & same family$^\dagger$ \\
P6 & bge-base (fp32) & bge-base (int8) & sanity only \\
P7 & all-MiniLM-L6-v2 & \todo{code encoder} & general$\to$code \\
P8 & \todo{code encoder A} & \todo{code encoder B} & code$\to$code \\
\bottomrule
\multicolumn{4}{@{}l}{\footnotesize $^\dagger$Overlaps the public registry of~\cite{embedflow}; used for calibration.}
\end{tabular}
\end{table*}

\subsection{Implementation}
Window, residual-miss and quality measurements use exact inner-product search, so
they are free of ANN error; HNSW~\cite{hnsw} with identical build parameters is used only
for P0 and for latency figures. BM25 uses \todo{Pyserini or bm25s} with Anserini's default
parameters ($k_1{=}0.9$, $b{=}0.4$). Adapter families and training-set size follow the
setup reported in~\cite{drift-adapter} \todo{copy exact families and $n$ from the paper}.
Hardware: \todo{GPU type, count}. Every encoding is computed once and cached, so all
strategies are replayed over identical vectors.

\subsection{Pipeline invariants}
Before any pair is interpreted, the replay pipeline must satisfy properties that hold
by construction: (i) P0 with exact search gives $\rho(K'){=}0$; (ii) $\rho(K')$ is
non-increasing in $K'$; (iii) re-ranking a window equal to the corpus gives $G{=}0$;
(iv) every trajectory with a completing background order ends at $M_T$ with $C{=}N$;
(v) S0 serves exactly $M_S$ until it switches. A violation is a pipeline bug; an
unexpected value on a real pair is a finding.

\subsection{Pre-registered hypotheses}
Fixed before running RQ1--RQ4; each with its falsification condition.
\begin{itemize}
\item \textbf{H1.} $G(K')$ at fixed $K'/N$ increases as compatibility $\kappa$
decreases. \emph{Falsified if} the Spearman correlation $r_s < 0.5$ across pairs.
\item \textbf{H2.} With the winner defined by the pre-registered rule of
Section~\ref{sec:frame} at the pre-registered $(\tau, r)$ setting, leave-one-pair-out prediction from $(\kappa, \rho(K'))$ beats the
majority-class baseline in mean regret. \emph{Falsified if} its mean regret $\ge$ that of
the majority baseline.
\item \textbf{H3 (exploratory).} At least one non-control pair has low compatibility
but low residual miss.
\item \textbf{H4.} The bounded-cost probe estimates $\rho(K')$ within $\pm$\todo{0.05}
of the full measurement at a \emph{total} encoding cost (probe windows, shared sample and
strata, deduplicated) of $\le$\todo{2\%} of $N$; the probe windows alone cost
$|Q_\text{probe}|\cdot K'$, about 1\% of $N$ for 20 queries at $K'/N = 5\cdot10^{-4}$. \emph{Falsified if} the error
exceeds the bound on more than one pair.
\item \textbf{H5.} The document signal of Section~\ref{sec:rq4} has AUC significantly
above 0.5 (bootstrap 95\% CI excludes 0.5) on at least half of the cross-family pairs,
and O-signal beats O-rand at equal budget. \emph{Otherwise} reported as a negative result.
\end{itemize}

% =====================================================================
\section{RQ1: Which Strategy Wins Where?}\label{sec:rq1}
% =====================================================================

\textbf{Protocol.} For every pair we encode the full corpus with both models (offline
evaluation only), replay each strategy--order combination along the migration, and
record $M_\sigma(C)$, $C_\varepsilon$, and the floor.

\begin{table}[t]
\centering
\caption{Per-pair measurements and the winner under the pre-registered rule (admissible: floor $\ge M_S-\delta$). \todo{fill}.}
\label{tab:map}
\small
\begin{tabular}{@{}lcccccl@{}}
\toprule
Pair & $\kappa$ & CKA & $\rho(K')$ & $G(K')$ & $C_{0.01}$ & Winner \\
\midrule
P1 & -- & -- & -- & -- & -- & -- \\
P2 & -- & -- & -- & -- & -- & -- \\
P3 & -- & -- & -- & -- & -- & -- \\
P4 & -- & -- & -- & -- & -- & -- \\
P5 & -- & -- & -- & -- & -- & -- \\
P6 & -- & -- & -- & -- & -- & -- \\
P7 & -- & -- & -- & -- & -- & -- \\
P8 & -- & -- & -- & -- & -- & -- \\
\bottomrule
\end{tabular}
\end{table}

\todo{Figure: scatter of pairs on ($\kappa$, $\rho(K')$), points labeled by winning
strategy, one panel per corpus.}

\todo{Figure: recall-vs-encodings trajectories for two contrasting pairs (e.g.\ P3 vs P2).}

\textbf{The empty-cell check.} \todo{Report whether any non-control pair falls in the
low-compatibility / low-residual-miss region (H3). If none does, state that window
re-ranking wins only where mapping also works, and discuss what that implies.}

\textbf{Text versus code.} \todo{Compare the position of code pairs on the map with text
pairs; test whether lexical nomination (S4) closes the gap on code, where identifiers
carry much of the relevance signal.}

% =====================================================================
\section{RQ2: Is the Winner Predictable?}\label{sec:rq2}
% =====================================================================

A map is a method only if its axes predict the winner on an unseen pair. We hold out
each pair in turn, fit a decision rule (a depth-2 tree and a nearest-neighbor rule) on
the remaining pairs' $(\kappa, \rho(K'))$, and predict the held-out winner. We compare
with a majority-class baseline and with each single axis.

\todo{Table: LOPO accuracy and regret (cost-to-$\varepsilon$ of predicted vs.\ true
winner) for: both axes, $\kappa$ only, $\rho$ only, majority.}

\note{With 8--10 pairs, report regret, not only accuracy: a wrong but near-tied
prediction costs little.}

% =====================================================================
\section{RQ3: A Bounded-Cost Probe}\label{sec:rq3}
% =====================================================================

Measuring $\rho(K')$ directly requires the native target top-$K$, i.e.\ encoding all of
$D$ with $f_t$ --- the very cost migration seeks to avoid. We estimate it instead.

\textbf{Intruder-count estimator.} For a probe query $q$, encode the window
$\win_q = \topk_{K'}^{s}(q)$ with $f_t$ and let $\tau_q$ be the $K$-th highest target
score in $\win_q$. A document outside the window can displace a window document from
the target top-$K$ only if its target score exceeds $\tau_q$. Define the intruder count
\[
I_q = \lvert\{d \notin \win_q : s_t(q,d) > \tau_q\}\rvert .
\]
$\min(I_q, K)$ is an \emph{upper bound} on the number of target top-$K$ documents missing
from the window, not its value: two intruders that beat only the weakest window
documents may displace just one of them. The exact count follows from intruder counts
at every window score. Let $w_1 \ge \dots \ge w_K$ be the top-$K$ target scores in
$\win_q$ and $I_q(w) = \lvert\{d\notin\win_q : s_t(q,d) > w\}\rvert$. The overall top-$K$
contains $j$ outside documents iff the $j$-th best outside score exceeds $w_{K-j+1}$, so
\[
\mathrm{miss}_q = \bigl\lvert\{\, j \in 1..K : I_q(w_{K-j+1}) \ge j \,\}\bigr\rvert .
\]
Drawing a uniform sample $S\subset D$ of size $m$ and encoding it with $f_t$ gives the
unbiased count estimate
\begin{gather*}
\hat I_q(w) = \frac{N}{m}\,\lvert\{d \in S\setminus \win_q : s_t(q,d) > w\}\rvert ,
\end{gather*}
which we plug into $\mathrm{miss}_q$; we also report the conservative bound
$\min(\hat I_q(w_K), K)/K$.
One sample serves every probe query, so the cost is
$\lvert\bigcup_q \win_q \cup S\rvert \le |Q_{\text{probe}}|\cdot K' + m$ encodings; all of them
count toward migration, so none are wasted. The window term dominates at scale: with
$N{=}10^6$ and $K'{=}500$, twenty probe queries already cost 1\% of $N$, so the number of
probe queries, not only $m$, is the budget to tune.

\textbf{The difficulty.} Intruders are rare tail events, so a uniform sample has high
variance when $I_q$ is small relative to $N$. We evaluate (i) the uniform estimator,
(ii) a stratified estimator that oversamples documents near $q$ in the \emph{source}
space (deeper source ranks $K' < r \le cK'$), with importance weights, and (iii) a
qrels-based variant that encodes only judged-relevant documents outside $\win_q$.
\note{Compare against EmbedFlow's finite-tail T2-v1 diagnostic, which avoids native
target rankings by design~\cite{embedflow}.}

\todo{Table: estimation error of $\hat\rho$ vs.\ full $\rho$ per pair, per estimator,
at encoding budgets of 0.1\%, 0.5\%, 1\% of $N$.}

\todo{Does the probe place each pair in the same map region as the full measurement,
and does LOPO prediction from probe estimates match RQ2?}

% =====================================================================
\section{RQ4: Document-Level Migration Ordering (Gated)}\label{sec:rq4}
% =====================================================================

\textbf{Signal.} Some documents are hot under the source model because of lexical
evidence (shared identifiers, rare terms); others are hot only because of the source
model's semantics. We hypothesize the latter are the ones whose exposure changes under
the target model. Let $\mathrm{kocc}_s(d)$ be the number of workload queries for which
$d\in\topk_K^{s}(q)$. The \emph{dense-only heat} of $d$ is
\[
h(d) = \bigl\lvert\{q : d\in\topk_K^{s}(q) \wedge d\notin\topk_{K_b}^{\mathrm{BM25}}(q)\}\bigr\rvert ,
\]
computable from the old index and a BM25 index, before any target encoding.

\textbf{Gate.} Documents with large source exposure mechanically show large absolute
exposure change, so the gate uses ratios: the signal is the dense-only share
$h(d)/\mathrm{kocc}_s(d)$ and the label is whether the relative change
$\lvert\mathrm{kocc}_t(d)-\mathrm{kocc}_s(d)\rvert/\mathrm{kocc}_s(d)$ exceeds its median, over
documents with $\mathrm{kocc}_s(d)>0$ (on the probe sample in deployment).
O-signal is enabled for a pair only if the bootstrap 95\% CI of the AUC excludes 0.5.

\textbf{Head-to-head.} At equal encoding budget we compare O-signal, O-rand, O-size,
and O-demand on the trajectory metrics, under ORCAS workloads (text) and the Zipf
sweep (code).

\todo{Table: AUC with CI per pair; gate decision.}
\todo{Figure: trajectory difference O-signal $-$ O-demand vs.\ Zipf $\alpha$.}

\note{Prior evidence is cautionary: in content-change maintenance, an exposure signal
added nothing over change history~\cite{cipbench}. If H5 fails, this section reports a
negative result with the measured AUCs.}

% =====================================================================
\section{Discussion}\label{sec:discussion}
% =====================================================================

\todo{Practitioner guidance distilled from the map: a short decision procedure
derived from RQ2/RQ3, with the regret it incurs.}

\todo{Where mapping and window strategies disagree, which one fails first and why
(hubness~\cite{hubness}, query--document asymmetry, instruction-prefixed encoders).}

\todo{Interaction with index-level reuse~\cite{dgm}: whether old-graph structure that
accelerates new-index construction also predicts window reliability.}

\textbf{Threats to validity.}
\begin{itemize}
\item \emph{Pair coverage:} \todo{N} pairs cannot represent all upgrades; winner prediction is
evidence about these pair types, not a universal law.
\item \emph{Subsampled corpora:} $\rho$ depends on $K'/N$; we report it as a fraction and
include full-scale points to check extrapolation.
\item \emph{Workloads:} ORCAS heat is real for text only; code skew is parametric.
\item \emph{Encoding cost as the cost unit} understates query-time latency for window
strategies; we report both.
\item \emph{Encoder correctness:} every model is encoded under an explicit contract (prompt,
pooling, padding side, EOS, revision) that is checked against the model card's worked
example and for batched-versus-single agreement; models that fail are excluded, not tuned.
\end{itemize}

% =====================================================================
\section{Related Work}\label{sec:related}
% =====================================================================

\textbf{Embedding model upgrades.} Drift-Adapter maps new queries into the legacy
space to defer re-encoding~\cite{drift-adapter}. EmbedFlow serves a target model over
source-index candidates while materializing target vectors, with an empirical
window diagnostic and a registry of measured pairs~\cite{embedflow}. Residual
reachability reuses old graph structure to accelerate new-index construction~\cite{dgm}.
Concurrent work studies exactness conditions and the economics of incomplete
migration~\cite{thaw-anon}. Our contribution is orthogonal: a common frame and a
cross-strategy, cross-domain map with predictive validation.

\textbf{Representation alignment.} Linear and orthogonal alignment of embedding
spaces~\cite{procrustes}, representation similarity~\cite{cka}, and relative
representations~\cite{relrep} motivate our compatibility axis.

\textbf{Dynamic vector indexes.} SPFresh, FreshDiskANN, and Quake update indexes as
vectors change~\cite{spfresh,freshdiskann,quake}; they assume the new vector exists,
which is the cost studied here. LEANN recomputes embeddings at query time to save
storage~\cite{leann}, a mechanism related to window re-ranking but with a different goal.

\textbf{Stale representations.} Corrector networks reduce re-encoding of stale
embeddings during retriever training~\cite{corrector}; CIPBench studies which content
changes require re-embedding under a fixed encoder~\cite{cipbench}.

% =====================================================================
\section{Conclusion}
% =====================================================================
\todo{Write after results.}

\section*{Artifacts}
\todo{Repository URL (anonymized): replay framework, cached encodings or scripts to
regenerate them, probe estimators, all per-pair tables.}

% =====================================================================
\begin{thebibliography}{99}
\bibitem{drift-adapter} H.~Vejendla, ``Drift-Adapter: A practical approach to near
zero-downtime embedding model upgrades in vector databases,'' in \emph{Proc. EMNLP},
2025, pp.~15938--15949.
\bibitem{embedflow} A.~Srivastav, ``EmbedFlow: Progressive embedding-model migration
over existing vector indexes,'' software, v0.8.0, 2026.
\url{https://github.com/arnsri33/embedflow}
\bibitem{dgm} \todo{authors}, ``Exploiting residual reachability for cross-model
migration of graph-based indexes in approximate nearest neighbor search,''
arXiv:2609.23036, 2026.
\bibitem{thaw-anon} Anonymous, ``\todo{Thaw title},'' under submission, 2026.
\bibitem{cipbench} N.~Zhao, J.~Zhang, T.~Zhong, H.~Wang, and B.~Wang, ``Which code
changes need re-embedding? CIPBench and the limits of multi-signal index
maintenance,'' \emph{ACM TOSEM} (submitted), 2026.
\bibitem{zilliz-migration} Zilliz, ``Migrate to a new embedding model,'' documentation,
2026.
\bibitem{msmarco} T.~Nguyen \emph{et al.}, ``MS MARCO: A human generated machine reading
comprehension dataset,'' arXiv:1611.09268, 2016.
\bibitem{orcas} N.~Craswell \emph{et al.}, ``ORCAS: 18 million clicked query-document
pairs for modeling search,'' in \emph{Proc. CIKM}, 2020.
\bibitem{codesearchnet} H.~Husain \emph{et al.}, ``CodeSearchNet challenge: Evaluating
the state of semantic code search,'' arXiv:1909.09436, 2019.
\bibitem{bm25} S.~Robertson and H.~Zaragoza, ``The probabilistic relevance framework:
BM25 and beyond,'' \emph{Found. Trends Inf. Retr.}, vol.~3, no.~4, 2009.
\bibitem{hnsw} Y.~Malkov and D.~Yashunin, ``Efficient and robust approximate nearest
neighbor search using hierarchical navigable small world graphs,'' \emph{IEEE TPAMI},
vol.~42, no.~4, 2018.
\bibitem{cka} S.~Kornblith, M.~Norouzi, H.~Lee, and G.~Hinton, ``Similarity of neural
network representations revisited,'' in \emph{Proc. ICML}, 2019.
\bibitem{procrustes} P.~Sch\"onemann, ``A generalized solution of the orthogonal
Procrustes problem,'' \emph{Psychometrika}, vol.~31, 1966.
\bibitem{relrep} L.~Moschella \emph{et al.}, ``Relative representations enable zero-shot
latent space communication,'' in \emph{Proc. ICLR}, 2023.
\bibitem{hubness} M.~Radovanovi\'c, A.~Nanopoulos, and M.~Ivanovi\'c, ``Hubs in space:
Popular nearest neighbors in high-dimensional data,'' \emph{JMLR}, vol.~11, 2010.
\bibitem{spfresh} Y.~Xu \emph{et al.}, ``SPFresh: Incremental in-place update for
billion-scale vector search,'' in \emph{Proc. SOSP}, 2023.
\bibitem{freshdiskann} A.~Singh \emph{et al.}, ``FreshDiskANN: A fast and accurate
graph-based ANN index for streaming similarity search,'' arXiv, 2021.
\bibitem{quake} J.~Mohoney \emph{et al.}, ``Quake: Adaptive indexing for vector search,''
in \emph{Proc. OSDI}, 2025.
\bibitem{leann} \todo{authors}, ``LEANN: A low-storage vector index,'' arXiv, 2025.
\bibitem{corrector} N.~Monath \emph{et al.}, ``A fresh take on stale embeddings:
Improving dense retriever training with corrector networks,'' in \emph{Proc. ICML},
2024.
\end{thebibliography}

\end{document}
